\chapter{Conclusiones}
\label{sec.cap9:conclusiones}

El análisis desarrollado a lo largo del presente informe permite concluir que la prevención y protección contra
incendios en instalaciones industriales y de ingeniería electrónica conforma un sistema integral de barreras múltiples
donde la termodinámica, la ingeniería eléctrica, el marco legal y la gestión del factor humano se encuentran
indisociablemente vinculados.

En primer término, la caracterización de los tipos de fuego (Clases A, B, C, D y K) demuestra que no existe una solución
única de extinción. La elección del agente extintor exige un conocimiento profundo del combustible involucrado y de los
riesgos asociados. En entornos donde predominan equipos electrónicos y tableros bajo tensión (Clase C), el empleo de
agentes tradicionales como el agua o el polvo químico seco triclase resulta contraproducente o destructivo: el agua
introduce el peligro letal de electrocución, mientras que el polvo químico ejerce una acción corrosiva e higroscópica que
inutiliza permanentemente los circuitos integrados y las pistas de cobre. Por tal motivo, los agentes gaseosos limpios
(fluorocetonas FK-5-1-12, gases halocarbonados y mezclas inertes tipo Inergen), el dióxido de carbono y los sistemas de
agua nebulizada a alta presión emergen como las alternativas tecnológicas superiores, capaces de suprimir las llamas sin
dejar residuos colaterales dañinos.

En segundo término, la física de la ignición de origen eléctrico —regida por el efecto Joule en contactos aflojados por
fatiga termomecánica y por la formación de arcos eléctricos en serie— evidencia la insuficiencia de las protecciones
termomagnéticas tradicionales para evitar incendios. La incorporación de dispositivos de detección de defecto de arco
(AFDD) bajo norma IEC 62606 y la adopción de inspecciones predictivas mediante termografía infrarroja constituyen
medidas de ingeniería indispensables para erradicar las fallas antes del encendido de los aislantes poliméricos.

En tercer lugar, el análisis toxicológico confirma que el monóxido de carbono y el cianuro de
hidrógeno actúan como agentes letales de asfixia química y celular en cuestión de minutos, provocando la desorientación y
el colapso muscular de las víctimas antes de que puedan visualizar las llamas. Esta realidad resalta la importancia de
los sistemas de detección ultra-temprana, tales como la aspiración de alta sensibilidad (ASD/VESDA) y las barreras
lineales, los cuales brindan una ventana temporal valiosa para evacuar las instalaciones antes de que se alcancen
concentraciones atmosféricas intolerables.

Asimismo, el marco regulatorio vigente en la República Argentina, sustentado en la Ley Nº 19.587, el Decreto 351/79 y la
reglamentación AEA 90364, provee un andamiaje analítico riguroso para la determinación de la Carga de Fuego, el cálculo de
las Unidades de Ancho de Salida y la verificación térmica de los conductores ($I_{cc}^2 t \le k^2 S^2$). No obstante, las
tragedias históricas analizadas en el país y en el exterior subrayan que los cálculos teóricos resultan estériles si no se
garantiza un mantenimiento auditado y redundante de las fuentes de energía y de las reservas hídricas contra incendios.

Finalmente, la efectividad del mejor diseño tecnológico queda supeditada a la capacitación continua de los recursos
humanos. El adiestramiento práctico periódico con fuego real en el uso de matafuegos mediante la técnica PASAR, la
claridad jerárquica de la brigada de emergencias, la disciplina de desplazamiento agachado bajo el humo y la ejecución
rigurosa de simulacros con recuento de personal en puntos de encuentro seguros representan las verdaderas garantías para
salvaguardar la vida humana, fin supremo de toda la disciplina de Higiene y Seguridad en el Trabajo.
